\section{总结与延伸}

\subsection{从 token relation 到 datapoint relation}

本讲最值得保留的统一图景是：attention 并不天然属于语言 token，它是“让一个对象依据一组其他对象更新表示”的计算模板。Transformer 在 tokens 之间学习关系；NPT 将每个 datapoint 变成对象，再通过 reshape 在 datapoints 与 attributes 两条轴之间交替推理。

\begin{importantbox}{十二条核心结论}
\begin{enumerate}
  \item Self-attention 令 query、key、value 来自同一集合，学习输入内部关系。
  \item Multi-head attention 提供多个投影子空间，但不保证每个 head 都可被人类命名。
  \item Teacher forcing 用 gold prefixes 并行训练所有 next-token positions。
  \item Causal mask 防止模型读取未来 targets，是 likelihood 正确性的边界。
  \item NPT 仍是 parametric neural network；“non-parametric”指 inference 显式依赖训练 datapoints。
  \item $(X,M)$ 把 observed evidence 与 masked queries 放进统一 dataset representation。
  \item ABD 在 rows 之间 attention，ABA 在每行 attributes 之间 attention。
  \item Row permutation equivariance 保证 datapoint 顺序不携带虚假语义。
  \item Feature masking 提供 regularization，target masking 训练 relational lookup。
  \item UCI average-rank evidence支持竞争力，不支持所有任务普遍胜出。
  \item Corruption 与 duplicate intervention 证明若干任务中跨 datapoint mechanism 确实被使用。
  \item Quadratic datapoint attention 是核心限制，未来需要 retrieval、sparsity 或 representative points。
\end{enumerate}
\end{importantbox}

\subsection{一个工程化判断框架}

若要在新任务中判断 NPT-style mechanism 是否值得采用，可以依次问五个问题：第一，其他 examples 在 inference 时是否真的提供额外信息；第二，是否能合法、安全、低延迟地访问这些 examples；第三，batch/retrieval set 是否覆盖关键 neighbors；第四，普通 tree、k-NN、retrieval baseline 是否已足够；第五，corruption 与 intervention 能否证明模型没有只靠自身 features 或 shortcut。

\begin{knowledgebox}{实践清单}
先实现强 per-row baseline；再加入 explicit retrieval；随后比较固定 k-NN、learned retriever 与 full datapoint attention；最后做 row corruption、neighbor removal、label intervention、batch-composition sensitivity 和 OOD evaluation。只有性能与机制证据同时改善，才应把收益归因于跨 datapoint reasoning。
\end{knowledgebox}

\subsection{拓展阅读}

建议按“机制--架构--验证”顺序阅读：先回到 \emph{Attention Is All You Need} 理解 causal training；再读 Kossen 等人的 NPT 论文第 2 节与附录掌握 tensor shapes 和 masking；随后对照 Deep Sets、Set Transformer、Neural Relational Inference 与 Neural Processes，比较 permutation symmetry、message passing、context conditioning 和 uncertainty assumptions。工程上还应继续追踪 sparse attention 与 retrieval-augmented modeling，因为它们直接面对 NPT 的全数据 quadratic bottleneck。

\teachervoice{讲者的收尾并未宣称问题已经解决，而是把 scaling 与更丰富任务留作开放方向。最稳妥的总结是：NPT 证明了“训练一个会读取数据集的预测算法”可以工作，并提供了少见的 mechanism-level experiments；它同时把检索覆盖、计算复杂度与数据泄漏变成不可回避的系统问题。}

\end{document}
